\chapter{Présentation}

\section{Ce qu'est \soft{}}
\soft{} est un lecteur de musique pour Android consacré aux musiques des ordinateurs, consoles et bornes d'arcade
des années 1980 à 2000. Ces musiques ne sont pas des enregistrements : ce sont des programmes et des données qui
pilotent une puce sonore. Pour les écouter fidèlement, \soft{} reproduit, par logiciel, le comportement de ces puces.
On parle d'\emph{émulation}.

\begin{enclair}
Un fichier MP3 contient le son lui-même. Un fichier de musique de Commodore~64 ou de Mega Drive contient plutôt une
partition et les réglages d'un instrument électronique précis. \soft{} joue cette partition sur une copie logicielle
de l'instrument d'origine.
\end{enclair}

\section{Origine}
\soft{} est dérivé de \emph{ZXTune}, lecteur libre créé pour les musiques du ZX Spectrum puis étendu à de nombreuses
plateformes. \soft{} en reprend le cœur de lecture et y ajoute :
\begin{itemize}
  \item des moteurs d'émulation choisis pour leur précision (par exemple ares pour la Super Nintendo, Ootake pour la
        PC Engine, reSIDfp pour le Commodore~64) ;
  \item une visualisation en temps réel de chaque voix et de l'état interne des puces, puce par puce ;
  \item une interface repensée : listes de lecture multiples, favoris, file d'attente, réglages en direct.
\end{itemize}

\section{Ce que \soft{} sait faire}
\begin{itemize}
  \item Lire \lhunaModulesCount{} formats de musique et ouvrir \lhunaContainersCount{} formats de conteneurs
        (archives, images de disquettes, compresseurs), voir la partie~\ref{part:compat}.
  \item Parcourir le stockage du téléphone et de grands catalogues en ligne.
  \item Afficher un oscilloscope par voix, la réponse du filtre du SID, et des tableaux de bord détaillés pour les
        puces SID, FM, SNES, PC Engine, NES, MT-32 et les puces PCM des bornes Sega.
  \item Fonctionner avec Android Auto et l'écran verrouillé.
\end{itemize}

\section{Configuration requise}
\begin{itemize}
  \item Android~4.1 ou plus récent pour la version complète.
  \item Pour l'oscilloscope : OpenGL ES, ou Vulkan quand l'appareil le permet (plus fluide).
  \item Espace de stockage : environ 16~Mo pour l'application, plus vos musiques.
\end{itemize}

\section{Historique}
\subsection{ZXTune}
ZXTune est né dans la communauté du ZX Spectrum, dont la scène musicale a produit des dizaines de milliers de
morceaux pour la puce AY-3-8910 avec des trackers comme Sound Tracker, Pro Tracker ou ASC Sound Master. Chacun de
ces logiciels enregistrait ses musiques dans son propre format, et les lecteurs de l'époque ne savaient en lire
qu'un seul. ZXTune, créé par Vitamin, s'est donné pour but de tous les lire, de les retrouver même cachés dans des
images de disquettes, des archives ou des programmes compressés, et de les rejouer fidèlement.

Le projet s'est ensuite ouvert aux autres plateformes en intégrant les meilleures bibliothèques libres existantes :
libsidplayfp pour le Commodore~64, Game Music Emu pour les consoles, libopenmpt et libxmp pour les modules de
tracker, vgmstream pour l'audio des jeux récents, libvgm pour les enregistrements VGM. Il existe en ligne de
commande, avec une interface graphique pour ordinateur et en application Android.

\subsection{ZXTune\_LHUNA}
\soft est une branche de l'application Android de ZXTune, centrée sur deux objectifs :
\begin{description}
  \item[La fidélité] choisir pour chaque puce le moteur d'émulation le plus juste disponible, et pouvoir comparer
        avec un moteur plus léger quand l'appareil peine ;
  \item[La lisibilité] montrer ce que fait chaque puce, voix par voix et image par image, avec des vues pensées
        pour la synthèse propre à chaque puce plutôt qu'une vue générique.
\end{description}
Tout ce qui est affiché est l'état réel de la puce émulée à l'instant présent : pas de valeurs moyennées ni
inventées. Le texte complète toujours le graphique, pour qu'un auditeur non technicien comprenne ce qu'il voit et
qu'un musicien retrouve exactement les valeurs de ses registres.
